\begin{abstract}

Machine learning-based spam detection faces persistent challenges from data scarcity, privacy constraints, and class imbalance, limiting development of robust classifiers. While Large Language Models (LLMs) offer promise for generating synthetic training data through neural text generation, fundamental questions remain about their effectiveness compared to traditional feature-space augmentation methods and across different classifier architectures. We present a rigorous evaluation framework employing systematic replication ($R = 20$ independent groups) and comprehensive statistical analysis to compare LLM-based text-space generation against SMOTE feature-space interpolation. Through 1,760 experimental configurations spanning three generation methods (GPT-4.1-mini, Claude-3.5-Haiku, SMOTE), three prompt strategies, two mixing approaches, eleven synthetic proportions (0\%-100\%), and two classifier architectures (SVM, Random Forest), we establish quantitative relationships between augmentation paradigms and classification performance. Our findings reveal striking architectural dependence: Random Forest demonstrates 11.8-fold greater robustness than SVM for LLM methods (2.1\% vs. 24.1\% degradation at 100\% synthetic), with fundamentally different degradation patterns (SVM: $R^2 = 0.990$ linear, Random Forest: $R^2 = 0.723$ non-linear). The SMOTE comparison crystallizes fundamental paradigm differences: feature-space interpolation exhibits extreme variance ($\pm$0.083, 4.2$\times$ higher than LLMs) with distribution-dependent outcomes ranging from +22.5\% improvement to -27.7\% degradation, while LLM text generation provides stable, predictable modest degradation ($\pm$0.020) with interpretable output. Contrary to expectations, prompt engineering intensity shows negligible impact (effect sizes $d < 0.1$), while distribution matching proves less critical than anticipated. The rigorous statistical framework—incorporating hypothesis testing with False Discovery Rate correction, Cohen's $d$ effect sizes, confidence intervals, and sensitivity analysis—transforms synthetic data evaluation from point estimation to confirmatory hypothesis testing with quantified uncertainty. These findings provide evidence-based guidance for operational deployment: prioritize ensemble classifiers, choose SMOTE for maximum Random Forest performance when pilot validation confirms benefits, select LLMs when stability and text interpretability are paramount, and focus optimization on model selection rather than prompt tuning. The systematic methodology establishes evaluation standards for synthetic cybersecurity data research while enabling practitioners to select between augmentation paradigms based on operational requirements and risk tolerance.

\end{abstract}
